\chapter{Rust pour la faible latence}\label{ch:ll:rust-for-low-latency}

Le jumeau Rust du décodeur de messages de ce livre décode un message en environ onze nanosecondes sur ce portable, à moins d'un cinquième
de l'original C++ d'un passage à l'autre. Les deux vérifient chaque longueur avant de lire, aucun n'alloue par message, et la
version Rust a été écrite sans un seul bloc \texttt{unsafe}. Ce que Rust change n'est pas la vitesse mais l'endroit où les erreurs
sont détectées~: une \omterm{def:ll:rust-for-low-latency:race}{course aux données} ou une référence pendante est une erreur de compilation, un indice hors bornes donne une panique
contrôlée au lieu d'une lecture silencieuse de la mémoire d'autrui. Ce chapitre examine Rust sur le chemin critique (hot path) de la
même manière que les chapitres~6 à~8 ont examiné C++~: quelles abstractions ne coûtent vraiment rien, où se trouvent les
vérifications et quand le compilateur les supprime, comment écrire du code \texttt{unsafe} qui reste sûr, comment maîtriser
l'allocation et les threads, et où le mécanisme asynchrone de Rust a sa place dans un système de trading et où il n'en a pas.

\section{Abstractions à coût nul, vérifiées}

\begin{definition}[Abstraction à coût nul]\label{def:ll:rust-for-low-latency:zero}
Une \emph{abstraction à coût nul}\index{abstraction à coût nul} se compile en le code qu'un programmeur aurait écrit à la main
sans elle~: la commodité de l'abstraction est résolue à la compilation et ne laisse aucune trace à l'exécution. La formule vient
du C++~: «~ce que vous n'utilisez pas, vous ne le payez pas~; et ce que vous utilisez, vous ne sauriez pas le coder mieux à la main~»
(Stroustrup).
\end{definition}

Les itérateurs, les fermetures, les génériques et \texttt{Option} sont conçus pour être à coût nul en Rust~: ils sont monomorphisés comme
les modèles C++ et intégrés en ligne. L'affirmation se vérifie de la même façon qu'au chapitre~7, dans l'assembleur. Le Rust sûr ajoute
une chose que le C++ n'a pas~: chaque indice de tranche est vérifié par rapport à la longueur de la tranche.

\begin{definition}[Élimination des vérifications de bornes]\label{def:ll:rust-for-low-latency:bce}
L'\emph{élimination des vérifications de bornes}\index{élimination des vérifications de bornes} est la suppression, par le compilateur,
d'un contrôle d'indice dont il peut prouver qu'il réussit toujours~: un compteur de boucle borné par la longueur propre de la tranche, un
indice déjà comparé, un itérateur qui ne produit que des positions valides.
\end{definition}

Le \cref{lst:ll:rust-for-low-latency:counted} est l'assembleur d'une boucle \texttt{for i in 0..v.len() \{ s += v[i] \}} compilée
par \texttt{rustc -O}~: il ne reste ni comparaison avec la longueur ni appel au gestionnaire de panique, et LLVM a vectorisé la
boucle quatre éléments à la fois. Le \cref{lst:ll:rust-for-low-latency:checked} lit un élément à un indice que le compilateur ne
connaît pas~: une comparaison, un branchement prédit, et le chemin de panique placé hors de la ligne. Lorsque la vérification ne peut
pas être supprimée, elle coûte peu, pourvu qu'elle soit prédite et que le chargement qu'elle protège soit la partie lente~: dans la
\cref{fig:ll:rust-for-low-latency:twins}, un gather de 65\,536 éléments à travers une table d'indices dispersés coûte la même
chose, à la marge d'erreur près, avec et sans vérification, dans les deux langages.

\omcode{code/low-latency/09-rust-for-low-latency/rust/asm/sum_counted.s}{1}{40}{Une boucle comptée sur une tranche, rustc 1.97 -O~: pas de vérification de bornes, vectorisée.}\label{lst:ll:rust-for-low-latency:counted}

\omcode{code/low-latency/09-rust-for-low-latency/rust/asm/get_checked.s}{1}{12}{Une lecture indicée à un indice inconnu~: une comparaison et un branchement vers le gestionnaire de panique.}\label{lst:ll:rust-for-low-latency:checked}

\begin{omfigure}
\begin{tikzpicture}
\begin{axis}[name=dec,width=4.6cm,height=5.4cm,ybar,bar width=12pt,ymin=0,symbolic x coords={decode},xtick=data,xticklabels={décodage},
  enlarge x limits=0.8,title={\small ns par message},tick label style={font=\footnotesize},label style={font=\small},
  grid=major,grid style={gray!25},table/col sep=comma]
\addplot[fill=omDef!55,draw=omDef] table[x=task,y=cpp]{figdata/low-latency/09-rust-for-low-latency/measured_decode.csv};
\addplot[fill=omThm!45,draw=omThm] table[x=task,y=rust]{figdata/low-latency/09-rust-for-low-latency/measured_decode.csv};
\end{axis}
\begin{axis}[at={(dec.east)},anchor=west,xshift=1.5cm,width=8.6cm,height=5.4cm,ybar,bar width=10pt,ymin=0,
  symbolic x coords={counted loop,gather checked,gather unchecked},xtick=data,enlarge x limits=0.25,
  xticklabels={{boucle\\comptée},{gather\\avec contrôle},{gather\\sans contrôle}},xticklabel style={align=center},
  title={\small ns par élément},tick label style={font=\footnotesize},label style={font=\small},grid=major,grid style={gray!25},
  legend columns=2,legend style={font=\footnotesize,at={(0.3,-0.33)},anchor=north,draw=none,fill=none,
  /tikz/every even column/.append style={column sep=8pt}},table/col sep=comma]
\addplot[fill=omDef!55,draw=omDef] table[x=task,y=cpp]{figdata/low-latency/09-rust-for-low-latency/measured_loops.csv};
\addplot[fill=omThm!45,draw=omThm] table[x=task,y=rust]{figdata/low-latency/09-rust-for-low-latency/measured_loops.csv};
\legend{C++ (g++ -O2),Rust (rustc -O)}
\end{axis}
\end{tikzpicture}

\omcaption{Les jumeaux C++ et Rust~: décodage de l'échantillon du livre~1 avec le décodeur de chaque langage (à gauche), et somme de
65\,536 entiers dans l'ordre, puis à travers une table d'indices dispersés avec et sans vérifications de bornes (à droite). Mesuré sur un
portable (Intel Core Ultra 7 155H) sous WSL2, sans cœurs isolés, épinglé sur un seul processeur. Données~: \texttt{bench\_rust.py}.}\label{fig:ll:rust-for-low-latency:twins}
\end{omfigure}

\section{Code unsafe et ses limites}

Les garanties de Rust valent pour le code sûr~; les blocs \texttt{unsafe} sont l'endroit où le programmeur reprend à son compte les
obligations de preuve du compilateur~: déréférencer des pointeurs bruts, appeler des fonctions C, contourner les vérifications de bornes,
implémenter des structures sans verrou (chapitre~11).

\begin{definition}[Abstraction saine]\label{def:ll:rust-for-low-latency:sound}
Une abstraction construite sur du code \texttt{unsafe} est \emph{saine}\index{abstraction saine} (sound) si aucune suite d'appels à son interface
sûre, aussi inhabituelle soit-elle, ne peut provoquer de \omterm{def:ll:cpp-for-latency-iii-the-compiler:ub}{comportement indéfini}~: les invariants sur lesquels s'appuie le code unsafe sont
établis et maintenus par l'abstraction elle-même, et non supposés chez ses appelants.
\end{definition}

La discipline consiste à garder la zone unsafe réduite et à consigner, à côté de chaque bloc, la raison pour laquelle il est sûr. Le
gather du \cref{lst:ll:rust-for-low-latency:sound} ne saute la vérification par élément qu'après que les indices ont été validés
une fois, et refuse une tranche dont la longueur n'est pas celle contre laquelle ils ont été validés~: quel que soit l'appel, il ne peut
pas lire hors bornes. Une fonction dont la sûreté dépend de son appelant est elle-même marquée \texttt{unsafe}, avec une section
\texttt{\# Safety} qui énonce le contrat~; \texttt{clippy} rejette une fonction publique qui déréférence un argument pointeur brut sans cela.
Les mêmes questions se posent pour le code C++, où chaque ligne est unsafe~; Rust les rend locales.

\omcode{code/low-latency/09-rust-for-low-latency/rust/src/lib.rs}{58}{72}{Une interface sûre au-dessus d'un gather sans contrôle.}\label{lst:ll:rust-for-low-latency:sound}

\section{Maîtriser l'allocation}

\begin{definition}[Allocateur global]\label{def:ll:rust-for-low-latency:alloc}
L'\emph{allocateur global}\index{allocateur global} d'un programme Rust est l'allocateur qui se cache derrière \texttt{Box}, \texttt{Vec},
\texttt{String} et toutes les collections standard~; un programme peut le remplacer par tout type implémentant \texttt{GlobalAlloc},
déclaré avec \texttt{\#[global\_allocator]}, par exemple un type qui compte les allocations dans les tests.
\end{definition}

Les types standard de Rust allouent à des endroits légèrement différents de ceux du C++, et les différences comptent sur un chemin critique. Une
\texttt{String} n'a pas de tampon intégré, de sorte qu'un symbole de quatre lettres alloue déjà~; un \texttt{Vec} réserve de la place pour
quatre petits éléments à sa première insertion, si bien que trois remplissages coûtent une allocation au lieu de trois~; \texttt{HashMap::entry}
déplace la clé dans la table au lieu de la copier. Mesurée avec \texttt{firm\_arena::CountingAlloc}, la version Rust du premier jet du
gestionnaire du chapitre~6 (\cref{lst:ll:rust-for-low-latency:naive}) allouait cinq fois par message là où la version C++ allouait neuf fois,
et la version corrigée n'alloue jamais. Le remède est le même qu'en C++~: tableaux de taille fixe, réservoirs d'indices, tableaux d'octets
pour les identifiants, pas de fermetures boxées sur le chemin critique, et un test qui vérifie le compte.

\omcode{code/low-latency/09-rust-for-low-latency/rust/src/lib.rs}{89}{102}{Le premier jet du chapitre~6 en Rust~: cinq allocations par message.}\label{lst:ll:rust-for-low-latency:naive}

\section{Threads, épinglage et ce que garantit le système de types}

\begin{definition}[Course aux données]\label{def:ll:rust-for-low-latency:race}
Une \emph{course aux données}\index{course aux données} se produit lorsque deux threads accèdent simultanément à la même zone mémoire, qu'au moins
un accès est une écriture et qu'au moins un n'est pas atomique. En C++, c'est un \omterm{def:ll:cpp-for-latency-iii-the-compiler:ub}{comportement indéfini}~; en Rust sûr, elle ne peut pas
s'écrire~: les traits \texttt{Send} et \texttt{Sync}, qui marquent les types pouvant être déplacés ou partagés entre threads, et la règle
selon laquelle une référence mutable est exclusive, sont vérifiés à la compilation.
\end{definition}

Le crate du chapitre contient un test de documentation marqué \texttt{compile\_fail}~: un thread qui incrémente un élément d'un
vecteur pendant que son propriétaire l'écrit aussi. Le test ne réussit que si le compilateur rejette le programme, ce qui est le cas,
car le thread lancé emprunte le vecteur de manière mutable alors qu'il est encore utilisé. Ce que le système de types ne fournit pas
c'est le placement~: la bibliothèque standard ne peut pas épingler un thread à un processeur. \texttt{firm.affinity} déclare les deux fonctions C dont il a besoin,
les appelle dans des blocs \texttt{unsafe} en énonçant la raison, et vérifie le résultat avec \texttt{sched\_getcpu}
(\cref{lst:ll:rust-for-low-latency:pin})~; son jumeau C++ fait de même avec \texttt{pthread\_setaffinity\_np}. Les deux lisent
le plan rôle-processeur qu'écrit \texttt{firm.coreplan} (chapitre~4).

\omcode{code/firm/affinity/rust/src/lib.rs}{36}{45}{Épinglage du thread appelant depuis Rust à travers la bibliothèque C.}\label{lst:ll:rust-for-low-latency:pin}

\section{Où l'asynchrone a sa place et où il n'en a pas}

\begin{definition}[Moteur d'exécution asynchrone]\label{def:ll:rust-for-low-latency:async}
Un \emph{moteur d'exécution asynchrone}\index{moteur d'exécution asynchrone} (async runtime) exécute de nombreuses tâches écrites sous forme de fonctions \texttt{async} sur quelques threads~: une tâche
qui doit attendre (une socket, une minuterie) rend la main à l'ordonnanceur du moteur, qui la reprend lorsque l'événement
survient, éventuellement sur un autre thread.
\end{definition}

Le code asynchrone est la manière dont Rust sert à moindre coût de nombreuses connexions lentes, et une société de trading en possède beaucoup~: les interfaces REST des plateformes, les flux WebSocket
de données de marché de plusieurs dizaines de plateformes de crypto-actifs (chapitre~17), les sessions de drop copy, les services internes. Le chemin critique est le
cas inverse~: un thread occupé par rôle, épinglé, qui tourne sur son entrée et ne cède jamais la main. Un ordonnanceur qui peut déplacer une tâche entre
threads, ou la retarder derrière une autre, ajoute précisément la \omterm{def:ll:where-latency-comes-from:tail}{gigue} que le chemin critique existe pour éviter. Les auteurs du moteur le
plus utilisé le disent eux-mêmes~: il est conçu pour des applications dont les tâches passent l'essentiel de leur temps à attendre
des entrées et sorties, et non pour accélérer le calcul. La répartition habituelle est un chemin critique synchrone et épinglé, alimenté et vidé par
des anneaux sans verrou (chapitre~12), entouré de services asynchrones pour tout ce qui peut attendre.

\section{Tutoriel~: le jumeau Rust, mesuré}\label{tut:ll:rust-for-low-latency:tut}

\begin{tutorial}
\textbf{Objectif.} Vérifier les abstractions de Rust dans l'assembleur, compter ses allocations, épingler un thread et comparer les jumeaux.
\textbf{État final~:} les deux listings d'assembleur, les comptes d'allocations et la \cref{fig:ll:rust-for-low-latency:twins}.
\begin{tutsteps}
\item \textbf{Assembleur.} \texttt{python ll\_rustasm.py} compile le crate avec \texttt{rustc -O --emit asm} et extrait
\texttt{sum\_counted} et \texttt{get\_checked}, exportées avec \texttt{\#[no\_mangle]} afin que leurs noms soient stables.
\item \textbf{Sûreté et course détectée à la compilation.} \texttt{cargo test} exécute les tests unitaires et le test \texttt{compile\_fail}
de documentation~; clippy s'exécute avec \texttt{-D warnings}.
\item \textbf{Allocations.} Le test installe \texttt{CountingAlloc} comme \omterm{def:ll:rust-for-low-latency:alloc}{allocateur global} et vérifie cinq allocations
par message pour le brouillon, et aucune pour le gestionnaire corrigé.
\item \textbf{Épinglage.} \texttt{firm\_affinity::apply(plan, role)} épingle un thread et vérifie son processeur.
\item \textbf{Mesurer} avec \texttt{python bench\_rust.py}~: il construit le benchmark du crate avec
\texttt{cargo build --release} et le jumeau C++ avec g++, puis exécute les deux épinglés.
\end{tutsteps}
\textbf{Pour aller plus loin.} Remplacer la recherche linéaire du gestionnaire corrigé par une petite table à adressage ouvert et garder le
nombre d'allocations à zéro~; compiler le jumeau C++ avec \texttt{-O3} et voir alors quel langage a la boucle comptée la plus rapide.
\end{tutorial}

\section{Construction~: placement des threads par rôle}\label{bld:ll:rust-for-low-latency:affinity}

\begin{build}
\bfield{Rôle} Chaque thread critique de la partie~IV est épinglé au processeur que le plan de placement attribue à son rôle, et le signale~;
les mesures du chapitre~26 consignent où chaque thread s'est exécuté.
\bfield{Interface} Rust \texttt{firm\_affinity}~: \texttt{read\_plan(text)}, \texttt{read\_plan\_file(path)},
\texttt{pin\_current(cpu)}, \texttt{current\_cpu()}, \texttt{allowed()}, \texttt{apply(plan, role)}. C++20
\texttt{firm::affinity}~: \texttt{read\_plan(path)}, \texttt{pin\_current}, \texttt{current\_cpu}, \texttt{name\_current},
\texttt{apply}. Python \texttt{firm\_coreplan.write\_plan} et \texttt{read\_plan} produisent et lisent le fichier de plan.
\bfield{Règles} Un thread est épinglé avant d'allouer de la mémoire ou de toucher à ses données~; \texttt{apply} échoue, plutôt que de continuer
sans épinglage, si le rôle manque, si le noyau refuse, ou si le thread n'est pas sur son processeur après coup~; chaque appel \texttt{unsafe}
porte son argument de sûreté.
\bfield{Tests d'acceptation} \texttt{code/firm/affinity/}~: le plan d'exemple lu dans les deux langages~; un thread épinglé à un processeur
autorisé et vérifié~; un rôle manquant refusé. \texttt{code/firm/coreplan/}~: l'aller-retour du fichier de plan.
\bfield{Pour aller plus loin} Donner au thread son nom pour \texttt{top} et \texttt{perf}~; refuser de démarrer lorsque le plan place deux rôles critiques
sur des threads frères.
\end{build}

\begin{omsources}
\item B.~Stroustrup, «~Abstraction and the C++ machine model~», 2005.
\item The Rustonomicon, «~Send and Sync~»~; The Rust Reference, «~Unsafety~».
\item Tokio, tutoriel, «~When not to use Tokio~».
\end{omsources}

\section{Exercices}

\begin{exercise}[$\star$]\label{exo:ll:rust-for-low-latency:1}
D'après la \cref{fig:ll:rust-for-low-latency:twins}, de quel facteur le décodeur Rust est-il plus lent ou plus rapide que celui en C++, et
que coûte cela à 5 millions de messages par seconde~?
\end{exercise}

\begin{exercise}[$\star$]\label{exo:ll:rust-for-low-latency:2}
Lesquelles de ces opérations allouent en Rust~: \texttt{\textquotedbl{}ESZ6\textquotedbl{}.to\_string()}~; \texttt{Vec::<u64>::new()}~; le premier
\texttt{push} dedans~; \texttt{[0u8; 24]}~; \texttt{Box::new(3)}~?
\end{exercise}

\begin{exercise}[$\star$]\label{exo:ll:rust-for-low-latency:3}
Pourquoi le gestionnaire brouillon en Rust alloue-t-il cinq fois là où celui en C++ alloue neuf fois~? Expliquer la différence
point par point.
\end{exercise}

\begin{exercise}[$\star\star$]\label{exo:ll:rust-for-low-latency:4}
Lire le \cref{lst:ll:rust-for-low-latency:counted}~: combien d'éléments la boucle sur le vecteur ajoute-t-elle à chaque itération, et qu'est-ce qui traite les
longueurs qui ne sont pas multiples de quatre~?
\end{exercise}

\begin{exercise}[$\star\star$]\label{exo:ll:rust-for-low-latency:5}
Pourquoi supprimer la vérification de bornes dans la collecte ne la rend-il pas plus rapide ici~? Quand le ferait-il~?
\end{exercise}

\begin{exercise}[$\star\star$]\label{exo:ll:rust-for-low-latency:6}
\texttt{CheckedIndex} reste-t-il sûr si \texttt{sum} ne compare pas \texttt{v.len()} avec \texttt{len}~? Donner une
séquence d'appels qui le casserait.
\end{exercise}

\begin{exercise}[$\star\star\star$]\label{exo:ll:rust-for-low-latency:7}
\emph{Programmation.} Écrire la boucle critique du décodeur Rust sur \texttt{decode(\&data).flatten()} avec un \texttt{match} explicite sur
le résultat à la place, compter les allocations avec \texttt{CountingAlloc}, et vérifier dans l'assembleur l'absence de paniques. Qu'est-ce qui change~?
\end{exercise}

\begin{exercise}[$\star\star\star$]\label{exo:ll:rust-for-low-latency:8}
\emph{Trouver la faille.} «~Notre gestionnaire de données de marché est une tâche asynchrone sur un moteur multi-thread afin qu'il passe à l'échelle avec le
nombre de cœurs.~»
\end{exercise}

\section{Problème~: Le portage}

\begin{problem}[{Problème du week-end --- réécriture du chemin de flux en Rust}]\label{pb:ll:rust-for-low-latency:1}
Une société envisage de porter son feed handler et son constructeur de carnet d'ordres de C++ vers Rust. L'équipe mesure les jumeaux
(\cref{fig:ll:rust-for-low-latency:twins}), compte les allocations et lit l'assembleur.

\medskip\noindent\textbf{Partie I --- Vitesse.}
\begin{enumerate}
\item Combien coûte chaque décodeur par message, et quel est le rapport~?
\item Pourquoi la boucle comptée en Rust est-elle plus rapide que celle en C++ avec les options utilisées~?
\item Que faudrait-il pour que la boucle C++ atteigne le même niveau~?
\item Que disent les mesures gather au sujet des vérifications de bornes sur cette charge de travail~?
\end{enumerate}

\medskip\noindent\textbf{Partie II --- Allocation.}
\begin{enumerate}[resume]
\item Combien d'allocations par message fait chaque brouillon, et chaque gestionnaire corrigé~?
\item Quel type Rust alloue là où le type C++ n'allouait pas~?
\item Comment le test impose-t-il zéro allocation en Rust~?
\item Que se passe-t-il si une dépendance alloue sur le chemin critique sans que l'équipe le sache~?
\end{enumerate}

\medskip\noindent\textbf{Partie III --- Sûreté.}
\begin{enumerate}[resume]
\item Quelle classe de bogues le portage supprime-t-il à la compilation~?
\item Où \texttt{unsafe} subsiste-t-il dans le code Rust de la société, et pourquoi~?
\item Que promet une \omterm{def:ll:rust-for-low-latency:sound}{abstraction saine} à ses appelants~?
\item Que signifie une panique sur le chemin critique en production, et comment faut-il la traiter~?
\end{enumerate}

\medskip\noindent\textbf{Partie IV --- La décision.}
\begin{enumerate}[resume]
\item Énoncer le \emph{résultat nommé}~: le rapport Rust sur C++ du temps par message pour le décodeur, et les allocations par
message de chaque brouillon et gestionnaire corrigé.
\item La différence de vitesse est-elle une raison de ne pas porter~?
\item Que porterait-on en premier, et pourquoi~?
\item Où l'async Rust a-t-il sa place dans les systèmes de la société~?
\item Comment les deux langages partagent-ils le plan de placement et les jeux de test de la société~?
\item Que perdrait l'équipe en portant le code, hormis la vitesse~?
\item Comment garderait-on les deux implémentations équivalentes pendant la transition~?
\item En une phrase~: que change Rust pour une équipe de latence faible~?
\end{enumerate}
\end{problem}

\section{Questions d'entretien}

\begin{interviewq}[$\star$ \iqroles{développeur}]\label{iq:ll:rust-for-low-latency:1}
Que signifie «~\omterm{def:ll:rust-for-low-latency:zero}{abstraction à coût nul}~», et comment vérifieriez-vous cette affirmation pour une chaîne d'itérateurs~?
\end{interviewq}

\begin{interviewq}[$\star\star$ \iqroles{développeur}]\label{iq:ll:rust-for-low-latency:2}
Les vérifications de bornes rendent-elles Rust lent~? Quand coûtent-elles, et comment les supprimeriez-vous en toute sûreté~?
\end{interviewq}

\begin{interviewq}[$\star\star$ \iqroles{développeur}]\label{iq:ll:rust-for-low-latency:3}
Qu'est-ce qui rend une abstraction \texttt{unsafe} \omterm{def:ll:rust-for-low-latency:sound}{saine}~?
\end{interviewq}

\begin{interviewq}[$\star\star$ \iqroles{développeur}]\label{iq:ll:rust-for-low-latency:4}
Comment Rust empêche-t-il les \omterm{def:ll:rust-for-low-latency:race}{courses aux données}, et que n'empêche-t-il pas~?
\end{interviewq}

\begin{interviewq}[$\star\star$ \iqroles{développeur}]\label{iq:ll:rust-for-low-latency:5}
Utiliseriez-vous l'async Rust sur un chemin critique de trading~? Pourquoi~?
\end{interviewq}

\begin{interviewq}[$\star\star\star$ \iqroles{développeur}]\label{iq:ll:rust-for-low-latency:6}
Comment prouveriez-vous qu'un composant Rust ne fait aucune allocation sur le tas après le démarrage, y compris dans ses dépendances~?
\end{interviewq}
